\section{Aplikacja --- dokumentacja użytkownika}
\label{sec:aplikacja}

Aplikacja jest zestawem narzędzi w~języku Python (3.10+), uruchamianych z~wiersza
poleceń, oraz pomocniczej aplikacji webowej zbudowanej w~bibliotece Streamlit.
Dwie ścieżki pracy uzupełniają się:

\begin{itemize}
    \item \textbf{Tryb CLI} --- generacja mapy, generacja populacji metodą ABC,
          walidacja i~wizualizacja.
    \item \textbf{Tryb web (Streamlit)} --- interaktywny edytor/generator mapy,
          podgląd, prostsze algorytmy (greedy/random) i~moduł porównawczy wyników.
\end{itemize}

\subsection{Wymagania i instalacja}

\paragraph{Wymagania:} Python $\ge 3.10$, biblioteki:
\texttt{numpy}, \texttt{matplotlib}, \texttt{pandas}, \texttt{streamlit}.
Pełna lista znajduje się w~pliku \texttt{requirements.txt}:

\begin{lstlisting}[language=bash, caption={Zawartość pliku \texttt{requirements.txt}}]
streamlit>=1.35.0
pandas>=2.0.0
matplotlib>=3.8
numpy>=1.26
scipy>=1.11
\end{lstlisting}

Instalacja w~środowisku wirtualnym:
\begin{lstlisting}[language=bash]
python -m venv .venv
# Windows:
.venv\Scripts\activate
# Linux/Mac:
source .venv/bin/activate

pip install -r requirements.txt
\end{lstlisting}

\subsection{Charakterystyka danych}
\label{sec:dane}

\subsubsection{Format mapy --- CLI (\texttt{map.json})}

Mapa generowana skryptem \texttt{map\_generator.py} zapisywana jest do pliku
JSON. Schemat:

\begin{lstlisting}[language=json, caption={Schemat \texttt{map.json}}]
{
  "weight": [[int, ...], ...],   // siatka H x W wag (int >= 1)
  "time":    int,                // budzet czasu T_max
  "budget":  int,                // budzet pieniezny B
  "attractions": [
    [row, col, value, cost, type],  // atrakcja: 5 intow
    ...
  ],
  "attraction_types": [
    [time, min, max],            // typ: (czas zwiedzania, min, max)
    ...
  ]
}
\end{lstlisting}

\paragraph{Uwaga.} Atrakcje są listami pięciu intów; \texttt{type} to
indeks (0-based) typu w~tablicy \texttt{attraction\_types}.

\subsubsection{Format mapy --- aplikacja web (scenarios/*.json)}

Aplikacja Streamlit posługuje się rozszerzonym formatem (m.in.\ z~nazwą,
nazwami typów, polami startu/końca jako słowniki):

\begin{lstlisting}[language=json, caption={Schemat \texttt{scenarios/*.json} (wycinek z \texttt{small.json})}]
{
  "name": "Small 8x8 - 4 atrakcje",
  "width": 8, "height": 8,
  "weights": [[3,2,4,...], ...],
  "attractions": [
    {"id": "a1", "x": 1, "y": 5, "value": 5, "type": "zabytek"},
    ...
  ],
  "attraction_types": ["zabytek", "przyroda"],
  "start": {"x": 0, "y": 0},
  "end":   {"x": 7, "y": 7},
  "budget": 30,
  "seed": 7
}
\end{lstlisting}

\subsubsection{Format populacji (\texttt{initial\_population.json})}

\begin{lstlisting}[language=json]
[
  {
    "id": 0,
    "path": [[r, c], [r, c], ...],
    "visited_attractions": [[r, c], ...],
    "used_attractions":    [[r, c], ...],
    "movement_cost": float,
    "attraction_cost": int,
    "total_time": float,
    "total_value": int,
    "type_counts": [int, ...],
    "budget_ok": bool, "time_ok": bool, "type_ok": bool,
    "feasible": bool
  },
  ...
]
\end{lstlisting}

\subsection{Uruchamianie --- tryb CLI}

\subsubsection{Krok 1. Wygenerowanie mapy}

Skrypt \texttt{map\_generator.py} po uruchomieniu wprost generuje mapę
$30 \times 30$ z~trzema typami atrakcji i~zapisuje wynik do \texttt{map.json}:

\begin{lstlisting}[language=bash]
python map_generator.py
\end{lstlisting}

Domyślne wartości (do edycji w~bloku \texttt{if \_\_name\_\_ == "\_\_main\_\_"}):
\begin{itemize}
    \item rozmiar: $30 \times 30$,
    \item zakres wag pól: $[7, 12]$,
    \item budżet czasowy: $800$, budżet pieniężny: $300$,
    \item trzy typy atrakcji z~różnymi parametrami $(\tau, m, M)$
          i~zakresami wartości oraz kosztów.
\end{itemize}

\subsubsection{Krok 2. Generowanie populacji metodą ABC}

\begin{lstlisting}[language=bash]
python initial_population.py \
    --population-size 50 \
    --iters 200 \
    --limit 60 \
    --out initial_population.json
\end{lstlisting}

\paragraph{Parametry skryptu:}
\begin{description}[leftmargin=2.6cm,style=nextline]
    \item[\texttt{-{}-map}] ścieżka do pliku mapy JSON (domyślnie \texttt{map.json}).
    \item[\texttt{-{}-out}] plik wyjściowy populacji (domyślnie \texttt{initial\_population.json}).
    \item[\texttt{-{}-population-size}] rozmiar populacji $N$ (domyślnie 50).
    \item[\texttt{-{}-iters}] liczba iteracji ABC (domyślnie 200).
    \item[\texttt{-{}-limit}] limit stagnacji $L$ uruchamiający fazę zwiadowczą (domyślnie 60).
    \item[\texttt{-{}-onlookers}] liczba prób onlookerów na iterację (domyślnie $= N$).
    \item[\texttt{-{}-seed}] ziarno generatora liczb losowych (opcjonalne).
\end{description}

\paragraph{Wyjście:} plik JSON z~populacją oraz wypisane na stdout
statystyki populacji (\cref{lst:stats}) i~,,najlepsza trasa''.

\begin{lstlisting}[caption={Przykład wyjścia statystyk populacji (\texttt{print\_population\_stats}).},label=lst:stats]
=== STATYSTYKI POPULACJI POCZATKOWEJ ===
Liczba tras: 50
Spelnia wszystko: 50/50
Spelnia budzet: 50/50
Spelnia limit czasu: 50/50
Spelnia ograniczenia typow: 50/50

Sredni calkowity czas: ...
Sredni koszt ruchu: ...
Sredni uzyty budzet: ...
Sredni a liczba atrakcji: ...
Srednia wartosc atrakcji: ...
Srednia dlugosc sciezki (liczba pol): ...

Srednia liczba atrakcji kazdego typu:
  Typ 0: ... (min=..., max=..., time=...)
  ...
\end{lstlisting}

\subsubsection{Krok 3. Walidacja wygenerowanej populacji}

Skrypt \texttt{examples/validate\_population.py} sprawdza, czy:
\begin{itemize}
    \item każda trasa ma flagę \texttt{feasible=true},
    \item żadna krawędź (nieskierowana) nie występuje w~trasie dwa razy.
\end{itemize}

\begin{lstlisting}[language=bash]
python examples/validate_population.py initial_population.json
\end{lstlisting}

Zwraca kod wyjścia 0 w~razie pomyślnej walidacji.

\subsubsection{Krok 4. Wizualizacja tras}

\begin{lstlisting}[language=bash]
python visualize_paths.py
\end{lstlisting}

Domyślnie funkcja \texttt{draw\_paths} czyta \texttt{map.json}
i~\texttt{initial\_population.json}. Można wybrać podzbiór tras po id:

\begin{lstlisting}[language=Python]
from visualize_paths import draw_paths
draw_paths(selected_ids=[0, 1, 5, 10])
\end{lstlisting}

Lub wskazać inny plik populacji:

\begin{lstlisting}[language=bash]
python -c "from visualize_paths import draw_paths; draw_paths(population_file='initial_population_abc.json')"
\end{lstlisting}

Wynik: okno Matplotlib z~wizualizacją (atrakcje jako czarne punkty,
ścieżki w~różnych kolorach, start jako zielony kwadrat, koniec jako
czerwony X).

\subsection{Uruchamianie --- aplikacja web (Streamlit)}

\begin{lstlisting}[language=bash]
streamlit run app.py
\end{lstlisting}

Aplikacja udostępnia dwa tryby (wybierane w~bocznym pasku):

\subsubsection{Tryb ,,Edytor / Generator''}

W~lewej kolumnie znajduje się formularz generatora mapy oraz edytor
parametrów (rozmiar, liczba atrakcji, ziarno, zakresy wartości i~wag,
typy atrakcji, rozkład wag --- \texttt{uniform}/\texttt{clusters},
budżet, punkty startu i~końca). Po wygenerowaniu można dalej ręcznie
edytować mapę (nazwę, budżet, listę atrakcji w~tabeli) oraz wczytywać
i~zapisywać scenariusze (\texttt{scenarios/*.json}).

W~prawej kolumnie widoczna jest mapa z~wagami w~kolorach. Można uruchomić
jeden z~dwóch baselineów:
\begin{itemize}
    \item \textbf{Greedy (zachłanny)} --- \texttt{solve\_greedy\_attractions},
    \item \textbf{Random walk} --- \texttt{solve\_random\_walk}.
\end{itemize}

Pod wynikiem prezentowane są metryki: koszt, budżet, wartość, liczba
odwiedzonych atrakcji oraz informacja, czy ścieżka mieści się w~budżecie.

\subsubsection{Tryb ,,Wyniki''}

Pozwala wczytać plik wyników o~strukturze:

\begin{lstlisting}[language=json]
{
  "map_file": "scenarios/small.json",
  "runs": [
    {
      "algo": "GA",
      "path": [[x, y], ...],
      "history": [fitness_0, fitness_1, ...],
      "fitness": liczba
    },
    ...
  ]
}
\end{lstlisting}

W~zakładkach prezentowane są: \emph{mapa z~nałożonymi ścieżkami},
\emph{wykres konwergencji} (na podstawie \texttt{history}) i~\emph{tabela
zbiorcza} (\texttt{algorithm}, fitness, koszt, wartość, liczba atrakcji,
długość ścieżki, flaga \emph{feasible}).

\subsection{Ustawianie parametrów instancji}

Parametry instancji ustawiane są w~zależności od~ścieżki użycia:

\begin{itemize}
    \item \textbf{CLI} --- bezpośrednio w~pliku \texttt{map\_generator.py}
          (zmienne \texttt{rows}, \texttt{cols}, \texttt{time\_limit},
          \texttt{budget}, wywołania \texttt{add\_attraction\_type}),
          a~następnie uruchomienie \texttt{python map\_generator.py}.
    \item \textbf{Web} --- formularz w~lewej kolumnie trybu
          ,,Edytor / Generator'' (rozmiar, liczba atrakcji, zakresy
          wartości i~wag, typy, dystrybucja, start/koniec, budżet, seed).
          Można też wczytać gotowy scenariusz z~\texttt{scenarios/} lub
          z~uploadu pliku.
\end{itemize}

\subsection{Ustawianie parametrów algorytmu (ABC)}

Parametry sterujące zachowaniem ABC są dostępne jako argumenty linii
poleceń (patrz wcześniejszy podrozdział). W~praktyce:
\begin{itemize}
    \item Zwiększenie \texttt{-{}-iters} i~\texttt{-{}-onlookers} daje lepszą
          eksploatację, ale wydłuża czas.
    \item Niższy \texttt{-{}-limit} sprzyja eksploracji (częstsze użycie
          zwiadowców), wyższy --- eksploatacji.
    \item \texttt{-{}-seed} zapewnia powtarzalność wyników.
\end{itemize}

W~kodzie funkcja \texttt{generate\_abc\_population} przyjmuje też dodatkowy
parametr \texttt{deduplicate} (domyślnie \texttt{True}), pilnujący, by
populacja końcowa nie zawierała duplikatów ścieżek.

\subsection{Wyniki i sposób ich interpretacji}

\paragraph{Statystyki agregowane.} Funkcja
\texttt{print\_population\_stats} wypisuje liczność rozwiązań spełniających
poszczególne ograniczenia (budżet, czas, typy, wszystko), średnie wartości
metryk i~rozkład liczności typów. Pomaga to ocenić ,,zdrowie'' populacji
(np.\ niski odsetek \texttt{feasible} może oznaczać zbyt rygorystyczne
ograniczenia instancji).

\paragraph{Najlepsza trasa.} Skrypt wyświetla cechy najlepszego rozwiązania
wg leksykograficznej krotki porządkowej (Value $\uparrow$, Time $\downarrow$,
movement cost $\downarrow$, attraction cost $\downarrow$).

\paragraph{Wizualizacja.} W~oknie \texttt{visualize\_paths.py}:
\begin{itemize}
    \item czarne punkty --- atrakcje (każda atrakcja niezależnie od typu),
    \item linie kolorowe --- poszczególne trasy z~populacji,
    \item zielony kwadrat --- start,
    \item czerwony X --- koniec.
\end{itemize}

Im więcej linii ,,konwerguje'' do podobnego kształtu, tym silniejsza była
selekcja w~ABC (mniejsza różnorodność końcowa). Jeśli linie są mocno
różne i~rozproszone, populacja jest zróżnicowana --- co bywa pożądane jako
materiał wejściowy dla dalszej optymalizacji.

\paragraph{Walidacja.}
Skrypt \texttt{examples/validate\_population.py} sygnalizuje:
\begin{itemize}
    \item \texttt{bad\_no\_repeated\_edges} --- liczba tras łamiących zakaz
          powtarzania krawędzi (oczekiwane: 0),
    \item \texttt{infeasible} --- liczba rozwiązań niedopuszczalnych
          (oczekiwane: 0).
\end{itemize}
Sumarycznie wypisuje status \texttt{OK} lub \texttt{NOT OK}.
